\section{هفته‌ی ۱۳: مدل‌های سلسله‌مراتبی}

\begin{center}\rule{0.5\linewidth}{0.5pt}\end{center}

\subsection{واژه‌های کلیدی}

واژه‌های مهم این درس با معادل فارسی‌شان (حتی اگر در درس‌های قبل آمده
باشند).

\begin{longtable}[]{@{}
  >{\raggedright\arraybackslash}p{(\linewidth - 4\tabcolsep) * \real{0.3333}}
  >{\raggedright\arraybackslash}p{(\linewidth - 4\tabcolsep) * \real{0.3333}}
  >{\raggedright\arraybackslash}p{(\linewidth - 4\tabcolsep) * \real{0.3333}}@{}}
\caption{\centering واژه‌های کلیدی}\tabularnewline
\toprule\noalign{}
\begin{minipage}[b]{\linewidth}\raggedright
\lr{English}
\end{minipage} & \begin{minipage}[b]{\linewidth}\raggedright
فارسی
\end{minipage} & \begin{minipage}[b]{\linewidth}\raggedright
توضیح کوتاه
\end{minipage} \\
\midrule\noalign{}
\endfirsthead
\toprule\noalign{}
\begin{minipage}[b]{\linewidth}\raggedright
\lr{English}
\end{minipage} & \begin{minipage}[b]{\linewidth}\raggedright
فارسی
\end{minipage} & \begin{minipage}[b]{\linewidth}\raggedright
توضیح کوتاه
\end{minipage} \\
\midrule\noalign{}
\endhead
\bottomrule\noalign{}
\endlastfoot
\lr{hierarchical model} & مدل سلسله‌مراتبی (چندسطحی) & پارامترهای فردی از
یک توزیع گروهی می‌آیند \\
\lr{parameter} & پارامتر & کمیت قابل تنظیم مدل که بین افراد یا شرایط فرق
می‌کند \\
\lr{estimate} / \lr{estimator} & برآورد / برآوردگر & مقدار به‌دست‌آمده
برای پارامتر / روش به دست آوردن آن \\
\lr{trial} & کوشش \lr{(trial)} & یک بار انجام تکلیف توسط شرکت‌کننده \\
\lr{prior} & توزیع پیشین \lr{(prior)} & باور درباره‌ی پارامتر پیش از دیدن
داده \\
\lr{likelihood} & درست‌نمایی & احتمال داده به‌عنوان تابعی از پارامترها \\
\lr{prior predictive check} & بررسی پیش‌بینی پیشین & شبیه‌سازی داده از
پیشین برای سنجش معقول بودن آن \\
\lr{posterior} & توزیع پسین \lr{(posterior)} & باور درباره‌ی پارامتر بعد
از دیدن داده \\
\lr{posterior predictive check (PPC)} & بررسی پیش‌بینی پسین & مقایسه‌ی
داده‌ی شبیه‌سازی‌شده از پسین با داده‌ی واقعی \\
\lr{HMC} / \lr{NUTS} & \lr{Hamiltonian Monte Carlo} / \lr{NUTS} &
\lr{MCMC} مبتنی بر گرادیان \\
\lr{divergent transition} & \lr{divergence} (گذار واگرا) & هشدار
\lr{NUTS} برای ناحیه‌ای که درست کاوش نشده \\
\lr{complete} / \lr{no} / \lr{partial pooling} & ادغام کامل / بدون ادغام
/ ادغام جزئی \lr{(pooling)} & یک پارامتر برای همه / جدا برای هر فرد /
سلسله‌مراتبی \\
\lr{shrinkage} & انقباض به سمت میانگین \lr{(shrinkage)} & کشیده شدن
برآوردهای پرنویز به سمت میانگین گروه \\
\lr{centered} / \lr{non-centered parameterization} & پارامتری‌سازی
\lr{centered} / \lr{non-centered} & دو شکل نوشتن مدل سلسله‌مراتبی \\
\lr{group-level parameter} & پارامتر سطح گروه & میانگین و پراکندگی
پارامترها در جمعیت \\
\lr{fit} / \lr{fitting} & برازش & تنظیم پارامترها تا مدل با داده جور
شود \\
\lr{hit} & اصابت \lr{(hit)} & پاسخ «بله» وقتی علامت وجود دارد \\
\lr{sensitivity (d′)} & حساسیت & فاصله‌ی توزیع علامت و نویز \\
\lr{gradient} & گرادیان & جهت و شیب تغییر تابع \\
\lr{simulation} & شبیه‌سازی & تولید داده‌ی مصنوعی با اجرای مدل \\
\lr{reparameterization} & بازپارامتری‌سازی & نوشتن پارامتر در مقیاس دیگر،
مثلاً \ensuremath{\lambda} = \lr{e\^{}\ensuremath{\eta}} \\
\lr{standard error (SE)} & خطای استاندارد & پراکندگی برآورد در تکرارهای
فرضی آزمایش \\
\lr{benchmark} & \lr{benchmark} (مجموعه‌آزمون استاندارد) & مجموعه سؤال
برای مقایسه‌ی مدل‌های هوش مصنوعی \\
\end{longtable}

\begin{center}\rule{0.5\linewidth}{0.5pt}\end{center}

\subsection{اهداف یادگیری}

در پایان این هفته، دانشجویان می‌توانند:

\begin{enumerate}
\def\labelenumi{\arabic{enumi}.}
\tightlist
\item
  به‌طور مفهومی \textbf{توضیح دهند} که \lr{Hamiltonian Monte Carlo} /
  \lr{NUTS} چه فرقی با \lr{Metropolis} گام تصادفی دارد و یک
  \lr{divergence} نشانه‌ی چیست.
\item
  یک مدل سلسله‌مراتبی (چندسطحی) را \textbf{صورت‌بندی کنند} و \lr{partial
  pooling} و \lr{shrinkage} را \textbf{توضیح دهند}، در مقایسه با
  رویکردهای \lr{complete pooling} و \lr{no pooling} هفته‌ی 5.
\item
  یک پارامترسازی \lr{non-centered} را \textbf{صورت‌بندی کنند} و
  استفاده‌اش را \textbf{توجیه کنند}.
\item
  پسین‌های سطح گروه و سطح فرد را \textbf{تفسیر کنند} و طبق چک‌لیست هفته‌ی
  12 \textbf{گزارش کنند}.
\end{enumerate}

\begin{center}\rule{0.5\linewidth}{0.5pt}\end{center}

\subsection{مقدمه}

در هفته‌ی 5 با یک دوراهی روبه‌رو شدیم. برازش \textbf{میانگین} همه‌ی
شرکت‌کننده‌ها پایدار است ولی می‌تواند مدل را تحریف کند (خطای میانگین‌گیری) و
تفاوت‌های فردی را نادیده می‌گیرد. برازش \textbf{هر شرکت‌کننده جداگانه}
تفاوت‌ها را رعایت می‌کند، ولی با \lr{trial}های کم برای هر نفر برآوردها
پرنویزند. بعضی‌شان حدی یا بی‌نهایت می‌شوند (شرکت‌کننده‌ای با 50/50 \lr{hit}
بدون تصحیح \(d' = \infty\) دارد). آماره‌های مرحله‌ی دوم روی این برآوردها
هم نادیده می‌گیرند که بعضی از بقیه خیلی نامطمئن‌ترند.

\textbf{مدل‌های سلسله‌مراتبی} این دوراهی را حل می‌کنند. فرض می‌کنند هر
شرکت‌کننده پارامتر خودش را دارد، \emph{و} این پارامترهای فردی از یک توزیع
سطح گروه می‌آیند که میانگین و پراکندگی‌اش هم برآورد می‌شود. آن‌وقت هر برآورد
فردی از بقیه نیرو قرض می‌گیرد: برآوردهای پرنویز و حدی دقیقاً به اندازه‌ای
که عدم قطعیتشان ایجاب می‌کند به سمت میانگین گروه کشیده \lr{(«shrink»)}
می‌شوند. پارامترهای سطح گروه به سؤال جمعیت («حساسیت معمول چقدر است و
آدم‌ها چقدر با هم فرق دارند؟») با در نظر گرفتن همه‌ی منابع عدم قطعیت جواب
می‌دهند.

نوشتن نمونه‌گیر برای چنین مدل‌هایی با دست خودش یک پروژه‌ی پژوهشی است.
نمونه‌گیرهای دست‌سازِ هفته‌های 11\lr{--}12 همان چیزی‌اند که این جعبه‌ی سیاه را
برایتان شفاف می‌کنند.

\paragraph{این موضوع کجا به کار می‌رود؟}

\begin{itemize}
\tightlist
\item
  \textbf{بیشتر مطالعه‌های امروزی مدل‌سازی شناختی با شرکت‌کننده‌های زیاد}
  مدل‌های سلسله‌مراتبی برازش می‌کنند: مدل‌های \lr{drift-diffusion}
  سلسله‌مراتبی، مدل‌های یادگیری تقویتی سلسله‌مراتبی در
  روان‌پزشکی محاسباتی، و \lr{SDT} سلسله‌مراتبی در پژوهش حافظه.
\item
  \textbf{مطالعه‌های بالینی و رشدی با جلسه‌های کوتاه:} کودکان و بیماران
  نمی‌توانند صدها \lr{trial} انجام دهند، پس برآوردهای فردی پرنویزند؛
  \lr{partial pooling} آن‌ها را قابل استفاده می‌کند.
\item
  \textbf{آموزش و علوم اجتماعی:} دانش‌آموزان درون کلاس‌ها درون مدرسه‌ها
  (مدل‌های چندسطحی) مثال کتاب‌درسی ساختار سلسله‌مراتبی است.
\item
  \textbf{تحلیل ورزشی و پزشکی:} برآورد مهارت یک بازیکن یا نرخ موفقیت یک
  بیمارستان از چند مشاهده‌ی کم با کمک گرفتن از گروه، همین ایده است.
\end{itemize}

\begin{center}\rule{0.5\linewidth}{0.5pt}\end{center}

\Needspace{11\baselineskip}

\subsection{جزوه‌ی درس}

\paragraph{به زبان ساده}

وقتی هر نفر فقط کمی داده به ما می‌دهد، برآوردهای فردی پرنویزند و بعضی‌شان
بی‌معنی از آب درمی‌آیند (نمره‌ی کامل یعنی \(d' = \infty\)). یک \textbf{مدل
سلسله‌مراتبی} می‌گوید: هر نفر پارامتر خودش را دارد، \emph{ولی} آدم‌ها به هم
شبیه‌اند، پس پارامترهایشان از یک توزیع گروهی مشترک می‌آیند که میانگین و
پراکندگی آن را هم برآورد می‌کنیم. نتیجه یک مصالحه‌ی معقول است: کسی که
داده‌ی کم یا حدی دارد به سمت میانگین گروه کشیده می‌شود
(\textbf{\lr{shrinkage}})، و کسی که داده‌ی زیادی دارد برآوردی نزدیک به
برآورد خودش نگه می‌دارد. شدت این کشش به این بستگی دارد که آدم‌ها چقدر با
هم فرق دارند، و این از روی داده یاد گرفته می‌شود. بخش آخر هفته یک
\textbf{گردش کار \lr{(workflow)}} قدم‌به‌قدم برای انجام امن این کار است.

\subsubsection{\texorpdfstring{2. \lr{Hamiltonian Monte Carlo} و
\lr{NUTS} (به‌طور
مفهومی)}{2. Hamiltonian Monte Carlo و NUTS (به‌طور مفهومی)}}

\lr{Metropolis} گام تصادفی گام‌های کورکورانه پیشنهاد می‌دهد و در ابعاد
بالا یا در امتداد \lr{ridge}های همبسته بیشترشان را هدر می‌دهد (هفته‌های
11\lr{--}12). \textbf{\lr{Hamiltonian Monte Carlo}} \lr{(HMC)} منفی لگاریتم
پسین را یک منظره و بردار فعلی پارامتر را یک توپ بدون اصطکاک در نظر
می‌گیرد. به توپ یک ضربه‌ی تصادفی (تکانه) بدهید، بگذارید مدتی با دنبال کردن
گرادیان منظره بغلتد، و جایی را که به آن می‌رسد پیشنهاد دهید. پیشنهادها
می‌توانند دور باشند و هنوز در نواحی پراحتمال، پس پذیرش بالا و خودهمبستگی
پایین است. \textbf{\lr{NUTS}} \lr{(No-U-Turn Sampler)} زمان غلتیدن را
خودکار تنظیم می‌کند.

\textbf{\lr{Divergence}ها.} غلتیدن به‌صورت عددی در گام‌های کوچک شبیه‌سازی
می‌شود. جایی که انحنای پسین ناگهان تغییر می‌کند (معمولاً در «قیف‌ها»، در
ادامه)، شبیه‌سازی نادقیق می‌شود و \lr{NUTS} یک \textbf{\lr{divergence}}
\lr{(divergent transition)} را علامت می‌زند. حتی چند \lr{divergence} یعنی
پسین آنجا به‌طور قابل اعتماد کاوش نمی‌شود. نادیده‌شان نگیرید؛
\lr{reparameterize} کنید.

\begin{quote}
\textbf{پیشرفته:} \lr{HMC} دینامیک هامیلتونی با انرژی
\(H(\boldsymbol\theta,\allowbreak \mathbf{p}) = -\ln\tilde p(\boldsymbol\theta) + \tfrac12\mathbf{p}^\top M^{-1}\mathbf{p}\)
را با انتگرال‌گیر \lr{leapfrog} شبیه‌سازی می‌کند. دینامیک دقیق \(H\) را
پایسته نگه می‌دارد، پس پیشنهادها همیشه پذیرفته می‌شدند. خطای \lr{leapfrog}
یک تصحیح \lr{Metropolis} را لازم می‌کند، و یک \lr{divergence} مسیری است
که خطای انرژی‌اش منفجر می‌شود.
\end{quote}

\subsubsection{3. مدل‌های سلسله‌مراتبی}

\textbf{ساختار (برای \lr{SDT} با \(J\) شرکت‌کننده):}

\[
\begin{aligned}
&\text{\lr{Group level:}} && \mu_d \sim \mathcal{N}(1, 1),\ \ \sigma_d \sim \mathcal{N}^{+}(0, 1),\ \ \mu_c \sim \mathcal{N}(0, 1),\ \ \sigma_c \sim \mathcal{N}^{+}(0, 1)\\
&\text{\lr{Individual level:}} && d'_j \sim \mathcal{N}(\mu_d, \sigma_d),\ \ c_j \sim \mathcal{N}(\mu_c, \sigma_c), \quad j = 1, \dots, J\\
&\text{\lr{Data:}} && \text{\lr{hits}}_j \sim \mathrm{Binomial}\big(n_s, \Phi(d'_j/2 - c_j)\big),\ \ \text{\lr{FAs}}_j \sim \mathrm{Binomial}\big(n_n, \Phi(-d'_j/2 - c_j)\big)
\end{aligned}
\]

\textbf{سه راه ترکیب اطلاعات (مقایسه با هفته‌ی 5):}

\begin{longtable}[]{@{}
  >{\raggedright\arraybackslash}p{(\linewidth - 4\tabcolsep) * \real{0.3333}}
  >{\raggedright\arraybackslash}p{(\linewidth - 4\tabcolsep) * \real{0.3333}}
  >{\raggedright\arraybackslash}p{(\linewidth - 4\tabcolsep) * \real{0.3333}}@{}}
\caption{\centering 3. مدل‌های سلسله‌مراتبی}\tabularnewline
\toprule\noalign{}
\begin{minipage}[b]{\linewidth}\raggedright
رویکرد
\end{minipage} & \begin{minipage}[b]{\linewidth}\raggedright
مدل
\end{minipage} & \begin{minipage}[b]{\linewidth}\raggedright
فرضش چیست
\end{minipage} \\
\midrule\noalign{}
\endfirsthead
\toprule\noalign{}
\begin{minipage}[b]{\linewidth}\raggedright
رویکرد
\end{minipage} & \begin{minipage}[b]{\linewidth}\raggedright
مدل
\end{minipage} & \begin{minipage}[b]{\linewidth}\raggedright
فرضش چیست
\end{minipage} \\
\midrule\noalign{}
\endhead
\bottomrule\noalign{}
\endlastfoot
\lr{complete pooling} & یک \(d'\) و یک \(c\) برای همه &
\(\sigma_d = \sigma_c = 0\): بدون تفاوت فردی \\
\lr{no pooling} & \(d'_j\) و \(c_j\) جداگانه با پیشین‌های ثابت و مستقل &
\(\sigma \to \infty\): شرکت‌کننده‌ها هیچ چیزی درباره‌ی هم به ما نمی‌گویند \\
\textbf{\lr{partial pooling}} & سلسله‌مراتبی؛ \(\sigma_d\) و \(\sigma_c\)
\emph{برآورد می‌شوند} & داده‌ها تصمیم می‌گیرند شرکت‌کننده‌ها چقدر شبیه
هم‌اند \\
\end{longtable}

\textbf{\lr{shrinkage.}} با یک تقریب نرمال، برآورد سلسله‌مراتبی
شرکت‌کننده‌ی \(j\) میانگین وزنی (با وزن دقت) برآورد خودش \(\hat d_j\) (با
خطای استاندارد \(\mathrm{SE}_j\)) و میانگین گروه است:

\[
\mathbb{E}[d'_j \mid \text{\lr{data}}] \approx \frac{\hat d_j/\mathrm{SE}_j^2 + \mu_d/\sigma_d^2}{1/\mathrm{SE}_j^2 + 1/\sigma_d^2}.
\]

شرکت‌کننده‌هایی با داده‌ی \textbf{پرنویزتر} (\(\mathrm{SE}_j\) بزرگ‌تر)
\textbf{بیشتر} \lr{shrink} می‌شوند. اگر آدم‌ها واقعاً خیلی با هم فرق داشته
باشند (\(\sigma_d\) بزرگ)، \lr{shrinkage} ضعیف است؛ اگر شبیه هم باشند،
قوی است. این از روی داده یاد گرفته می‌شود.

\textbf{چرا \lr{shrinkage} خوب است، نه تحریف:} در کل شرکت‌کننده‌ها،
برآوردهای \lr{shrink}شده به‌طور میانگین \emph{به مقادیر واقعی نزدیک‌ترند}
تا برآوردهای فردی خام. این یک نتیجه‌ی معروف در آمار است (پارادوکس
\lr{Stein).} برآوردهای خام حدی تا حدی از روی شانس حدی‌اند، و
\lr{shrinkage} این بازگشت به میانگین را جبران می‌کند.

\textbf{سؤال‌های سطح گروه} مستقیماً با پسین‌های \(\mu_d\) (حساسیت معمول) و
\(\sigma_d\) (تفاوت‌های فردی) جواب داده می‌شوند، و مقایسه‌ی گروه‌ها با
افزودن اثرهای گروه یا شرط به سطح گروه انجام می‌شود.

\paragraph{خواندن ریاضی}

\begin{itemize}
\tightlist
\item
  \textbf{دو سطح.} سطح گروه: \(\mu_d\) (حساسیت معمول) و \(\sigma_d\)
  (آدم‌ها چقدر با هم فرق دارند). سطح فرد:
  \(d'_j \sim \mathcal{N}(\mu_d,\allowbreak \sigma_d)\)، یعنی حساسیت نفر \(j\) از
  توزیع گروه کشیده می‌شود. سطح داده: \lr{hit}ها و \lr{false alarm}های نفر
  \(j\) فقط به \(d'_j\) و \(c_j\) همان نفر بستگی دارند.
\item
  \(\mathcal{N}^{+}(0, 1)\): توزیع «نیم‌نرمال» \lr{(half-normal)}، یک
  منحنی نرمال که در 0 تا شده و فقط مقدارهای مثبت در آن مانده‌اند (انحراف
  معیار نمی‌تواند منفی باشد).
\item
  \textbf{فرمول \lr{shrinkage}} یک میانگین وزنی با دقت است (مثل هفته‌ی
  12): برآورد خود فرد \(\hat d_j\) دقت \(1/\mathrm{SE}_j^2\) دارد و
  توزیع گروه دقت \(1/\sigma_d^2\). مثال حل‌شده: \(\mathrm{SE} = 0.40\)
  دقت 6.25 می‌دهد و گروه 4، پس برآورد
  \((6.25 \times 2.4 + 4 \times 1.3)/10.25 = 1.97\) است: مقدار خام 2.4،
  40\% راه را به سمت میانگین گروه یعنی 1.3 می‌رود.
\item
  اگر آدم‌ها واقعاً خیلی با هم فرق داشته باشند (\(\sigma_d\) بزرگ)، دقت
  گروه کم و \lr{shrinkage} ضعیف است. اگر شبیه هم باشند، \lr{shrinkage}
  قوی است.
\end{itemize}

\subsubsection{\texorpdfstring{4. پارامترسازی \lr{centered} در برابر
\lr{non-centered}}{4. پارامترسازی centered در برابر non-centered}}

نوشتن مستقیم \(d'_j \sim \mathcal{N}(\mu_d,\allowbreak \sigma_d)\) (شکل
\textbf{\lr{centered}}) در پسین یک «قیف» می‌سازد: وقتی \(\sigma_d\) کوچک
است همه‌ی \(d'_j\)ها باید به هم فشرده شوند و وقتی بزرگ است می‌توانند پخش
شوند. \lr{NUTS} با این هندسه‌ی متغیر به زحمت می‌افتد و \lr{divergence}
تولید می‌کند، به‌خصوص با داده‌ی کم برای هر شرکت‌کننده. شکل
\textbf{\lr{non-centered}} انحراف‌های استانداردشده را نمونه‌گیری می‌کند و
\(d'_j\) را به‌طور قطعی می‌سازد:

\[
z_j \sim \mathcal{N}(0, 1), \qquad d'_j = \mu_d + \sigma_d z_j.
\]

همان مدل را توصیف می‌کند ولی هندسه‌ی خیلی دوستانه‌تری دارد. وقتی داده برای
هر شرکت‌کننده کم است، به‌طور پیش‌فرض از آن استفاده کنید.

\paragraph{خواندن ریاضی}

\(z_j \sim \mathcal{N}(0, 1)\) و \(d'_j = \mu_d + \sigma_d z_j\): به جای
نمونه‌گیری مستقیم \(d'_j\)، یک نمره‌ی استانداردشده‌ی «این فرد چقدر غیرعادی
است» یعنی \(z_j\) (بر حسب انحراف معیار گروه) را نمونه بگیر و \(d'_j\) را
از روی آن بساز. هر دو نسخه دقیقاً یک مدل را توصیف می‌کنند. تفاوت فقط در
چیزی است که نمونه‌گیر باید کاوش کند: در شکل \lr{centered}، وقتی
\(\sigma_d\) کوچک است، محدوده‌ی مجاز هر \(d'_j\) تقریباً به هیچ می‌رسد (یک
قیف)؛ در شکل \lr{non-centered}، \(z_j\)ها همیشه روی یک مقیاس‌اند، پس
اندازه‌ی گام نمونه‌گیر همه‌جا مناسب است.

\subsubsection{5. گردش کار بیزی}

\begin{enumerate}
\def\labelenumi{\arabic{enumi}.}
\tightlist
\item
  مدل و پیشین‌ها را (با توجیه) \textbf{تعریف کنید}.
\item
  \textbf{بررسی پیش‌بینی پیشین} (هفته‌ی 9): آیا پیشین داده‌های معقولی نتیجه
  می‌دهد؟
\item
  با چند زنجیره \textbf{برازش کنید}.
\item
  \textbf{تشخیص دهید:} \(\hat R\)، \lr{ESS}، \lr{divergence}ها،
  \lr{trace plot}ها (هفته‌ی 12).
\item
  \textbf{بررسی پیش‌بینی پسین} (هفته‌ی 10): آیا مدل برازش‌شده می‌تواند
  داده‌ها را بازتولید کند؟
\item
  \textbf{تفسیر و گزارش کنید.}
\end{enumerate}

\paragraph{برای یک مهندس}

\begin{itemize}
\tightlist
\item
  \lr{shrinkage} به سمت میانگین گروه همان پدیده‌ی تخمین‌گر
  \textbf{\lr{James\lr{--}Stein}} است که وقتی سه واحد یا بیشتر داریم، در
  مجموع خطای مربعی از تخمین‌های \lr{ML} جداگانه بهتر است.
\item
  \textbf{\lr{HMC/NUTS}} دینامیک همیلتونی (یک ذره‌ی بدون اصطکاک روی سطح
  منفی لگاریتم \lr{posterior)} را با یک انتگرال‌گیر عددی \lr{(leapfrog)}
  شبیه‌سازی می‌کند، مثل شبیه‌سازی‌های فیزیکی. یک \textbf{\lr{divergence}}
  یعنی انتگرال‌گیر ناپایدار شده چون اندازه‌ی گام برای انحنای محلی خیلی
  بزرگ است؛ مثل یک معادله‌ی دیفرانسیل سخت \lr{(stiff)} که با گام خیلی
  بزرگ انتگرال‌گیری شود.
\item
  \textbf{پارامتری‌سازی \lr{non-centered}} یک تغییر متغیر است که متغیرها
  را به مقیاس واحد نرمال می‌کند، یعنی حالت \lr{(conditioning)} مسئله را
  بهتر می‌کند؛ شبیه سفید کردن یا نرمال کردن ورودی‌ها پیش از یک الگوریتم
  تطبیقی.
\end{itemize}

\paragraph{برای مدرس}

\begin{itemize}
\tightlist
\item
  \textbf{\lr{shrinkage} را با یک مثال ورزشی یا امتحانی معرفی کنید.}
  دانشجویی که 3 از 3 سؤال را درست جواب داده احتمالاً «دانشجوی 100\%»
  نیست؛ دانشجویی با 95 از 100 احتمالاً نزدیک 95\% است. از کلاس بپرسید
  اطلاعات فردی و کلاسی را چطور ترکیب می‌کنند: خودشان \lr{partial pooling}
  را اختراع می‌کنند.
\item
  \textbf{نمودار \lr{shrinkage} بخش \lr{D lab} را نشان دهید:} برآوردهای
  خام، برآوردهای سلسله‌مراتبی و میانگین گروه روی یک شکل، با پیکان. ایده
  را ماندگار می‌کند.
\item
  \textbf{ریاضی \lr{HMC} را عمیق درس ندهید.} یک تصویر (یک توپ که روی
  منفی لگاریتم \lr{posterior} با یک ضربه‌ی تصادفی می‌غلتد) کافی است؛ مهارت
  عملی مهم، خواندن شاخص‌های تشخیصی است \lr{(divergence}، \(\hat R\)،
  \lr{ESS).}
\item
  \textbf{سؤال‌های رایج.}

  \begin{itemize}
  \tightlist
  \item
    \emph{«آیا \lr{shrinkage} نتایج را سوگیر نمی‌کند؟»} کمی، ولی خطای کل
    را کم می‌کند، پس برآوردها به‌طور میانگین به حقیقت نزدیک‌ترند.
  \item
    \emph{«آیا مدل سلسله‌مراتبی همیشه بهتر است؟»} وقتی افراد داده‌ی کافی
    دارند، نتایج شبیه برازش‌های جداگانه است؛ بیشترین سود وقتی است که برای
    هر نفر داده کم است.
  \item
    \emph{\lr{«divergence} برای نتایجم یعنی چه؟»} یعنی بخشی از
    \lr{posterior} به‌طور قابل اعتماد کاوش نشده. \lr{reparameterize} کنید
    و دوباره اجرا کنید تا هیچ \lr{divergence}ای نماند.
  \end{itemize}
\end{itemize}

\begin{center}\rule{0.5\linewidth}{0.5pt}\end{center}

\subsection{مثال حل‌شده}

\textbf{\lr{shrinkage} با دست.} پسین گروه \(\mu_d \approx 1.3\) و
\(\sigma_d \approx 0.5\) دارد (دقت گروه \(1/0.5^2 = 4\)). دو شرکت‌کننده
هر دو برآورد خام \(\hat d_j = 2.4\) دارند.

\begin{longtable}[]{@{}
  >{\raggedright\arraybackslash}p{(\linewidth - 8\tabcolsep) * \real{0.2000}}
  >{\raggedright\arraybackslash}p{(\linewidth - 8\tabcolsep) * \real{0.2000}}
  >{\raggedright\arraybackslash}p{(\linewidth - 8\tabcolsep) * \real{0.2000}}
  >{\raggedright\arraybackslash}p{(\linewidth - 8\tabcolsep) * \real{0.2000}}
  >{\raggedright\arraybackslash}p{(\linewidth - 8\tabcolsep) * \real{0.2000}}@{}}
\caption{\centering مثال حل‌شده}\tabularnewline
\toprule\noalign{}
\begin{minipage}[b]{\linewidth}\raggedright
شرکت‌کننده
\end{minipage} & \begin{minipage}[b]{\linewidth}\raggedright
\(\mathrm{SE}_j\)
\end{minipage} & \begin{minipage}[b]{\linewidth}\raggedright
دقت فردی \(1/\mathrm{SE}_j^2\)
\end{minipage} & \begin{minipage}[b]{\linewidth}\raggedright
برآورد سلسله‌مراتبی
\end{minipage} & \begin{minipage}[b]{\linewidth}\raggedright
انحراف معیار پسین
\end{minipage} \\
\midrule\noalign{}
\endfirsthead
\toprule\noalign{}
\begin{minipage}[b]{\linewidth}\raggedright
شرکت‌کننده
\end{minipage} & \begin{minipage}[b]{\linewidth}\raggedright
\(\mathrm{SE}_j\)
\end{minipage} & \begin{minipage}[b]{\linewidth}\raggedright
دقت فردی \(1/\mathrm{SE}_j^2\)
\end{minipage} & \begin{minipage}[b]{\linewidth}\raggedright
برآورد سلسله‌مراتبی
\end{minipage} & \begin{minipage}[b]{\linewidth}\raggedright
انحراف معیار پسین
\end{minipage} \\
\midrule\noalign{}
\endhead
\bottomrule\noalign{}
\endlastfoot
\lr{P1 (trial} کم) & 0.40 & 6.25 &
\((6.25 \times 2.4 + 4 \times 1.3)/(6.25 + 4) = 20.2/10.25 = \mathbf{1.97}\)
& \(1/\sqrt{10.25} = 0.31\) \\
\lr{P2 (trial} زیاد) & 0.20 & 25 &
\((25 \times 2.4 + 4 \times 1.3)/(25 + 4) = 65.2/29 = \mathbf{2.25}\) &
\(1/\sqrt{29} = 0.19\) \\
\end{longtable}

\textbf{تفسیر.} هر دو شرکت‌کننده \emph{به نظر} به یک اندازه حساس‌اند، ولی
برآورد \lr{P1} بر داده‌ی کمتری تکیه دارد، پس مدل کمتر به آن اعتماد می‌کند
و آن را 0.43 به سمت میانگین گروه جابه‌جا می‌کند. \lr{P2} فقط 0.15 جابه‌جا
می‌شود. هر دو انحراف معیار پسین از \lr{SE}های خام کوچک‌ترند، چون توزیع
گروه اطلاعات اضافه می‌کند. (در مدل کامل، \(\mu_d\) و \(\sigma_d\) خودشان
نامطمئن‌اند و هم‌زمان برآورد می‌شوند؛ این محاسبه برای توضیح آن‌ها را ثابت
می‌گیرد.)

\begin{center}\rule{0.5\linewidth}{0.5pt}\end{center}

\Needspace{11\baselineskip}

\subsection{\texorpdfstring{\lr{lab} (2 ساعت)}{lab (2 ساعت)}}

\subsubsection{اهداف}

\begin{enumerate}
\def\labelenumi{\arabic{enumi}.}
\tightlist
\item
  شبیه‌سازی یک آزمایش با 12 شرکت‌کننده و محاسبه‌ی برآوردهای فردی \lr{(no
  pooling).}
\item
  نمونه‌گیری با \lr{NUTS}؛ بررسی \(\hat R\)، \lr{ESS} و
  \lr{divergence}ها.
\item
  نمایش \lr{shrinkage} و مقایسه‌ی برآوردها با مقادیر واقعی.
\item
  اجرای یک بررسی پیش‌بینی پسین و تفسیر پارامترهای سطح گروه.
\end{enumerate}

\textbf{زمان‌بندی:} بخش \lr{A} 10 دقیقه · بخش \lr{B} 20 دقیقه · بخش
\lr{C} 25 دقیقه · بخش \lr{D} 30 دقیقه · بخش \lr{E} 20 دقیقه · بخش \lr{F}
15 دقیقه.

\begin{quote}
بخش \lr{C} را زود شروع کنید.
\end{quote}

\subsubsection{\texorpdfstring{بخش \lr{B: }بررسی پیش‌بینی پیشین}{بخش B: بررسی پیش‌بینی پیشین}}

پیشین دامنه‌ی وسیعی از نرخ‌های \lr{hit} و \lr{false alarm} را نتیجه می‌دهد
که بیشترشان بالای قطرند. حدود 30\% شرکت‌کننده‌های شبیه‌سازی‌شده شمارش
\lr{hit} یا \lr{false alarm} دقیقاً 0 یا 50 دارند. این‌ها بیشتر
شرکت‌کننده‌هایی‌اند که با حساسیت خیلی بالا و سوگیری قوی شبیه‌سازی شده‌اند و
\lr{false alarm} صفر تولید می‌کنند. امتحانش کنید و ببینید آیا پسین تغییر
می‌کند.

\subsubsection{\texorpdfstring{بخش \lr{D: shrinkage} و استنباط سطح گروه
(30 دقیقه)}{بخش D: shrinkage و استنباط سطح گروه (30 دقیقه)}}

یکی را که برآورد خامش حدی است پیدا کنید و توصیف کنید مدل سلسله‌مراتبی با
آن چه می‌کند. 2. آیا برآورد سلسله‌مراتبی به‌طور میانگین به حقیقت نزدیک‌تر
است؟ آیا برای \emph{هر} شرکت‌کننده نزدیک‌تر است؟ چرا این لازم نیست؟ 3.
\(\mu_d\) و \(\sigma_d\) درباره‌ی این جمعیت چه می‌گویند؟ \(\sigma_d\) را
با انحراف معیاری که برای شبیه‌سازی داده استفاده شد (0.5) مقایسه کنید و
تفاوت را توضیح دهید (راهنمایی: 12 شرکت‌کننده).

\subsubsection{\texorpdfstring{بخش \lr{F:} گزارش (15
دقیقه)}{بخش F: گزارش (15 دقیقه)}}

یک پاراگراف روش‌ها و نتایج طبق چک‌لیست هفته‌ی 12 بنویسید: مدل و پیشین‌ها (با
کلمات و نمادها)، نمونه‌گیر و تنظیمات، تشخیص‌ها (بیشترین \(\hat R\)، کمترین
\lr{ESS}، \lr{divergence}ها)، نتیجه‌ی \lr{PPC}، و نتیجه‌گیری‌های سطح گروه
(\(\mu_d\)، \(\sigma_d\)، \(\mu_c\) با بازه‌های 95\%؛
\(P(\mu_c > 0\mid\text{\lr{data}})\)).

\subsubsection{تمرین‌ها (به ترتیب سختی)}

برآورد \(d'\) آن را با \(\mu_d\) مدل سلسله‌مراتبی، و بررسی پیش‌بینی پسینش
را با بخش \lr{E} مقایسه کنید.

\lr{divergence}ها، \lr{ESS} و زمان اجرا را مقایسه کنید. بعد داده را به
10 + 10 \lr{trial} برای هر شرکت‌کننده کم کنید و دوباره مقایسه کنید.

\textbf{تمرین 3 \ensuremath{\star}\ensuremath{\star} (تعداد \lr{trial} نابرابر).} به نیمی از شرکت‌کننده‌ها
20 + 20 \lr{trial} و به نیمی دیگر 200 + 200 \lr{trial} بدهید (داده‌ی جدید
شبیه‌سازی کنید). نشان دهید \lr{shrinkage}، مثل مثال حل‌شده، به تعداد
\lr{trial}ها بستگی دارد.

تشخیص‌های \lr{Pareto-}\(k\) را بررسی کنید و بگویید برای پیش‌بینی
\lr{leave-one-participant-out} چه معنایی دارند.

\textbf{تمرین 5 \ensuremath{\star}\ensuremath{\star}\ensuremath{\star} (پیشرفته: اثر گروه).} دو گروه 12 نفره شبیه‌سازی کنید
که حساسیت میانگینشان 0.4 فرق دارد. مدل را با یک اثر گروه \(\beta\) روی
\(\mu_d\) گسترش دهید
(\(d'_j = \mu_d + \beta\cdot\text{\lr{group}}_j + \sigma_d z_j\))، \(\beta\)
را برآورد کنید و \(P(\beta > 0\mid\text{\lr{data}})\) را گزارش کنید. با یک
آزمون \lr{t} روی برآوردهای فردی \(d'\) مقایسه کنید.

\subsubsection{\texorpdfstring{\lr{lab check-off}}{lab check-off}}

نشان دهید: (1) خروجی تشخیص‌ها (بدون \lr{divergence}، \(\hat R\)،
\lr{ESS)}، (2) نمودار \lr{shrinkage} با مقایسه‌ی \lr{RMSE}، و (3) نمودار
\lr{PPC} همراه با پاراگراف بخش \lr{F.}

\textbf{گزارش پیشرفت پروژه (این هفته):} یک گزارش وضعیت یک‌صفحه‌ای تحویل
دهید (چه کارهایی انجام شده، یک نمودار برازش یا شبیه‌سازی اولیه، مسائل
باز). نمره ندارد؛ بازخورد کتبی ظرف یک هفته داده می‌شود.

\begin{center}\rule{0.5\linewidth}{0.5pt}\end{center}

\subsection{\texorpdfstring{ارتباط با هوش مصنوعی و
\lr{LLM}ها}{ارتباط با هوش مصنوعی و LLMها}}

\textbf{\lr{partial pooling} برای \lr{benchmark}های هوش مصنوعی.}
\lr{benchmark}ها به دسته‌های تکلیف تقسیم می‌شوند، و بعضی دسته‌ها فقط چند
آیتم دارند. گزارش دقت خام برای هر دسته عددهای پرنویز و گاهی حدی می‌دهد،
درست مثل برازش‌های فردی \lr{SDT} با \lr{trial}های کم. یک مدل سلسله‌مراتبی
دقت دسته‌ها را نمونه‌هایی از یک توزیع مشترک در نظر می‌گیرد و دسته‌های کوچک
را به سمت میانگین کلی \lr{shrink} می‌کند.

دسته‌ی 8آیتمی (خام 0.88، واقعی 0.63) به‌شدت پایین کشیده می‌شود، و برآوردهای
\lr{shrink}شده در مجموع حدود دو برابر دقیق‌ترند \lr{(RMSE} برابر 0.066 در
برابر 0.125).

جوابتان را به فرمول \lr{shrinkage} در جزوه‌ی درس این هفته ربط دهید.
میانگین‌های پسین را با برآوردهای \lr{empirical-Bayes} مقایسه کنید و
تشخیص‌ها را بررسی کنید. - \ensuremath{\star}\ensuremath{\star} بحث: یک جدول رتبه‌بندی «دسته‌ای که مدل \lr{X}
در آن ضعیف‌ترین است» را برجسته می‌کند، که یک دسته‌ی 5آیتمی است. به
نویسندگان چه می‌گویید؟

\begin{center}\rule{0.5\linewidth}{0.5pt}\end{center}

\subsection{اشتباهات رایج و برداشت‌های نادرست}

\begin{enumerate}
\def\labelenumi{\arabic{enumi}.}
\tightlist
\item
  \textbf{\lr{«shrinkage} نتایج را سوگیر می‌کند.»} کمی سوگیری را با کاهش
  زیاد واریانس معاوضه می‌کند و به‌طور میانگین برآوردهای دقیق‌تری می‌دهد.
  اصولی است، نه یک تصحیح موقت.
\item
  \textbf{نادیده گرفتن \lr{divergence}ها.} حتی چندتایشان نشانه‌ی کاوش
  غیرقابل اعتماد است.
\item
  \textbf{پارامترسازی \lr{centered} با داده‌ی کم.} قیف و \lr{divergence}
  می‌سازد. برای مدل‌های سلسله‌مراتبی با \lr{trial}های کم برای هر شرکت‌کننده،
  به‌طور پیش‌فرض از \lr{non-centered} استفاده کنید.
\item
  \textbf{پیشین‌های مبهم برای انحراف معیارهای گروه.} پیشین‌های خیلی پهن
  روی \(\sigma\) می‌توانند پیش‌بینی‌های پیشین نامعقول و مشکلات نمونه‌گیری
  بسازند. پیشین‌های نیم‌نرمال با مقیاس معقول معمولاً بهترند.
\item
  \textbf{گزارش فقط برآوردهای فردی.} پارامترهای سطح گروه اغلب به سؤال
  پژوهشی واقعی جواب می‌دهند.
\end{enumerate}

\begin{center}\rule{0.5\linewidth}{0.5pt}\end{center}

\subsection{خودآزمایی}

\textbf{پاسخ} \lr{complete pooling} (همه یکسان‌اند، \(\sigma = 0\))،
\lr{no pooling} (شرکت‌کننده‌ها بی‌ارتباط‌اند، پس هیچ اطلاعاتی به اشتراک
گذاشته نمی‌شود)، و \lr{partial pooling/}سلسله‌مراتبی (شرکت‌کننده‌ها از یک
توزیع مشترک می‌آیند که پراکندگی‌اش از داده برآورد می‌شود).

\textbf{2.} دو شرکت‌کننده برآورد خام \(d'\) یکسانی دارند. چرا ممکن است
مدل سلسله‌مراتبی یکی را خیلی بیشتر \lr{shrink} کند؟ \textbf{پاسخ}
\lr{shrinkage} به دقت هر برآورد فردی بستگی دارد. شرکت‌کننده‌ای که
\lr{trial} کمتری دارد \lr{(SE} بزرگ‌تر) اطلاعات کمتری دارد و قوی‌تر به سمت
میانگین گروه کشیده می‌شود.

\textbf{3.} چرا پارامترسازی \lr{non-centered} اغلب ترجیح دارد؟
\textbf{پاسخ} شکل \lr{centered} پسینی قیفی‌شکل می‌سازد (پارامترهای فردی
وقتی انحراف معیار گروه کوچک است باید به هم فشرده شوند) که \lr{NUTS} بد
کاوشش می‌کند و \lr{divergence} می‌دهد. نمونه‌گیری \(z_j\) استانداردشده و
محاسبه‌ی \(d'_j = \mu + \sigma z_j\) هندسه‌ای می‌دهد که کاوشش خیلی آسان‌تر
است.

\textbf{4.} \lr{divergence} چیست و باید با آن چه کرد؟ \textbf{پاسخ}
مسیری در \lr{NUTS} که شبیه‌سازی عددی‌اش شکسته است، معمولاً در ناحیه‌ای با
انحنای زیاد؛ نشان می‌دهد بخشی از پسین به‌طور قابل اعتماد کاوش نمی‌شود.

\textbf{5.} چرا \lr{HMC} معمولاً در ابعاد بالا خیلی بهتر از
\lr{Metropolis} گام تصادفی مخلوط می‌شود؟ \textbf{پاسخ} از اطلاعات گرادیان
برای حرکت‌های بلند و جهت‌دار که شکل پسین را دنبال می‌کنند استفاده می‌کند، پس
پیشنهادهای دور هم هنوز احتمال پذیرش بالایی دارند. پیشنهادهای گام تصادفی،
مگر اینکه خیلی ریز باشند، تقریباً همیشه رد می‌شوند.

\textbf{پاسخ} به ترتیب \(\sigma_d\) (انحراف معیار سطح گروه) و \(\mu_d\)
(میانگین سطح گروه).

\textbf{7.} مراحل گردش کار بیزی را فهرست کنید. \textbf{پاسخ} تعریف مدل و
پیشین‌ها \ensuremath{\leftarrow} بررسی پیش‌بینی پیشین \ensuremath{\leftarrow} برازش با چند زنجیره \ensuremath{\leftarrow} تشخیص (\(\hat R\)،
\lr{ESS}، \lr{divergence}ها، \lr{trace}ها) \ensuremath{\leftarrow} بررسی پیش‌بینی پسین \ensuremath{\leftarrow} در
صورت نیاز مقایسه‌ی مدل‌ها \ensuremath{\leftarrow} تفسیر و گزارش.

\begin{center}\rule{0.5\linewidth}{0.5pt}\end{center}

\subsection{تکلیف}

\begin{itemize}
\tightlist
\item
  \textbf{\lr{A5 (MCMC} از صفر)} در ابتدای درس هفته‌ی 14 تحویل داده
  می‌شود.
\item
  \textbf{گزارش پیشرفت پروژه} این هفته (بالا را ببینید).
\item
  تکلیف نمره‌دار جدیدی نداریم: هفته‌های باقی‌مانده برای پروژه‌ی پایانی است.
\end{itemize}

\begin{center}\rule{0.5\linewidth}{0.5pt}\end{center}

\subsection{مطالعه‌ی بیشتر}

\begin{itemize}
\tightlist
\item
  \textbf{\lr{Lee} \lr{\&} \lr{Wagenmakers} (2014)، فصل 11.} نظریه‌ی تشخیص
  علامت، از جمله همان مدل \lr{SDT} سلسله‌مراتبی‌ای که امروز صورت‌بندی
  کردید. \emph{(منبع اصلی.)}
\item
  \textbf{\lr{Kruschke} (2014)، فصل 9.} مدل‌های سلسله‌مراتبی: \lr{partial
  pooling} و \lr{shrinkage}، قدم‌به‌قدم.
\item
  \textbf{\lr{Farrell} \lr{\&} \lr{Lewandowsky} (2018)، فصل 9.} مدل‌سازی
  چندسطحی یا سلسله‌مراتبی داده‌های شناختی.
\end{itemize}

\emph{پیش‌نمایش هفته‌ی 14:} با جعبه‌ابزار کامل، گشتی در مدل‌هایی از حوزه‌های
مختلف علوم شناختی می‌زنیم؛ با دسته‌بندی (مدل \lr{exemplar} یعنی \lr{GCM)}
و سازوکار سرعت\lr{--}دقت مدل \lr{drift-diffusion} شروع می‌کنیم.
